% Section 3 — Models used
\section{Models Used}
\label{sec:models}
Each subsection states what the method assumes, how it was configured in this
study, and what comparison it supports. Textbook derivations are omitted;
implementation truth is in appendix C.

\subsection{Naive}
\label{sec:m-naive}
Naive repeats the last observed value across the 28-day horizon, using only
history preceding each origin. It is the reference floor: a method that cannot
improve on repeating yesterday adds no value. Its accuracy is modest (M5 MAE
1.18; Store 16.29), yet it wins 6 of 27 Store inventory policies, an early
indication that accuracy and cost can diverge.

\subsection{Seasonal Naive}
\label{sec:m-snaive}
Seasonal Naive repeats the same weekday from the preceding week, the cheapest
use of retail's weekly rhythm (figure 7). The implementation repeats
only the last seven observed values; an earlier version that reached beyond
the seasonal window was identified and corrected during hardening. It serves
as the seasonal reference run in every phase. Despite middling accuracy, it is
the third-lowest-cost Store inventory input (2,128.14), which suggests weekly
structure stocks shelves effectively even when point accuracy is unremarkable.

\subsection{Moving Average}
\label{sec:m-ma}
The moving average forecasts the mean of the last $w$ days, a low-pass filter
on noisy demand and the standard practitioner baseline: no estimated dynamics
to break. The global window was selected on validation (M5: $w=14$ at MAE
0.9206, against 0.9236 for $w=7$ and 0.9282 for $w=28$; Store: $w=7$ at 7.4057,
against 7.5113 and 7.9132), with seed 42. It is the simple method the
remaining rungs must justify themselves against. It nearly leads M5
forecasting (MASE 1.346 against LSTM's 1.316) and records the lowest Store
inventory cost (2,084.50).

\subsection{Simple Exponential Smoothing}
\label{sec:m-ses}
SES tracks a slowly moving demand level with higher weight on recent days
\citep{hyndman2018}; it is the minimal model with memory, estimating level
without trend or seasonality. The transparent implementation uses
validation-selected $\alpha=0.1$. SES performs consistently well without
leading any forecast board: second-lowest cost on M5 inventory (159.31) and
on Store inventory (2,113.82), and winner of 2 M5 sensitivity policies at
short lead with high stockout penalty.

\subsection{Double Exponential Smoothing}
\label{sec:m-des}
DES (Holt's method) extends SES with a trend component
($\alpha=\beta=0.2$, grid-selected), testing whether explicit trend modeling
helps. It does not help here. DES carries roughly double the inventory of the
leading methods (average stock 8.41 on M5 and 76.07 on Store) while recording
mid-range or worse costs, indicating that trend extrapolation on noisy retail
data accumulates stock without commensurate accuracy. Its 5 Store sensitivity
wins reflect Store leadership fragmentation rather than DES strength.

\subsection{Triple Exponential Smoothing}
\label{sec:m-tes}
TES (additive Holt-Winters, $0.3/0.3/0.3$) adds a weekly seasonal component to
DES, completing the classical toolkit in one model. It records the weakest
results in the study: M5 MASE 2.60 with cost 346.68, Store MASE 2.11 with cost
3,419.88, and the highest error in every archetype. The finding is that
combined trend-seasonality structure, as implemented here, overfits the
available signal.

\subsection{ARIMA}
\label{sec:m-arima}
ARIMA models the next value as a linear combination of recent values and
recent shocks after differencing toward stationarity
\citep{box1970,hyndman2018}, providing the statistical treatment of
autocorrelation that smoothing lacks. A fixed ARIMA(1,1,0) was adopted after
a validation sanity check; per-series order search was rejected as
overfitting-prone and computationally out of scope. Series with more than 85\%
zeros skip estimation and fall back to Naive rather than fitting noise. ARIMA
is the linear-memory representative: mid-range accuracy on both datasets (M5
MASE 1.51; Store 1.45) with a marked feasibility asymmetry discussed in
section 4.13 (22.65\% pre-check fallback on sparse M5 demand against
zero fallbacks on dense Store demand).

\subsection{SARIMA}
\label{sec:m-sarima}
SARIMA adds seasonal autoregression and moving-average terms, testing whether
explicit weekly autocorrelation improves on SES-style seasonality for dense
data. The configuration is SARIMA(1,1,0)(0,1,1)$_7$, Store-only, estimated on
the full 500 series by 8 origins (4,000 estimations, 4,000 successful, 2,382
seconds total). A prior 100-series subset is retained as an exploratory record
but is not the headline result. SARIMA was not estimated on M5: seasonal
differencing of predominantly zero series lacks justification, and the
Croston rung already addresses that setting. It is the second-ranked Store
forecaster (MASE 1.055) but ranks fifth on Store inventory cost (2,282.31).

\subsection{LSTM}
\label{sec:m-lstm}
The LSTM is a recurrent network with gating that carries signal across long,
noisy sequences \citep{hochreiter1997}. Estimated globally (one network per
dataset, 28-day lookback, per-series history-only scaling, chronological
per-series validation rebuilt during hardening to remove overlapping-window
leakage, seed 42), it tests whether pooled cross-series learning outperforms
per-series statistics. Within the evaluated model set it records the lowest
forecast error on both datasets and the lowest M5 inventory cost, but ranks
fourth on Store inventory cost; section 4.12 examines this split in
detail.

\begin{figure}[H]
\centering
\includegraphics[width=0.85\textwidth]{02_lstm_cell_gates.png}
\caption{The LSTM cell: gates control what is retained, discarded, and
emitted at each step, allowing information to persist across weeks of noisy
demand. Teaching figure from notebook 08.}
\label{fig:lstmcell}
\end{figure}

\subsection{Croston}
\label{sec:m-croston}
Croston's method \citep{croston1972} separates intermittent demand into two
sequences (nonzero demand size and interval between nonzero demands) and
smooths each independently, avoiding fitting zeros as signal. Estimated
M5-only with validation-selected $\alpha=0.1$; the 28-day forecast is flat at
the estimated rate. Accuracy is respectable (MAE 1.08) and its M5 service
level is the highest recorded (93.9\%), but its cost (174.81) exceeds its
bias-corrected variants below.

\subsection{SBA}
\label{sec:m-sba}
The Syntetos-Boylan approximation \citep{syntetos2005} multiplies the Croston
rate by $(1-\alpha/2)$ to correct its over-forecast bias. Estimated M5-only
with $\alpha=0.5$. It is the lowest-cost member of the family (160.66): a
small algebraic correction with measurable inventory value.

\subsection{TSB}
\label{sec:m-tsb}
The Teunter-Syntetos-Babai method \citep{teunter2011} replaces Croston's
interval smoothing with direct updating of demand-occurrence probability,
responding faster when demand patterns shift. Estimated M5-only at
$(\alpha,\beta)=(0.1,0.3)$. It records the family's lowest forecast error
(MAE 1.010, MASE 1.369) but the family's second-lowest cost (165.13): accuracy
and cost rank the family differently.

\subsection{Model-ladder comparison}
\label{sec:ladder-compare}
\Cref{tab:taxonomy} compresses the ladder: twelve methods on five rungs
across two datasets, with the two deliberate scope exceptions (SARIMA
Store-only, Croston family M5-only). The ladder orders complexity; whether
value follows is resolved empirically in section 4.

\begin{table}[H]
\centering
\caption{Model taxonomy: purpose, scope, and comparative role. Full configurations in appendix C.}
\label{tab:taxonomy}
\begin{tabular}{llll}
\toprule
Model & Rung & Scope & Comparative role \\
\midrule
Naive & Baseline & Both & Reference floor \\
Seasonal Naive & Baseline & Both & Weekly-rhythm reference \\
Moving Average & Baseline & Both & Practitioner benchmark \\
SES & Smoothing & Both & Level memory \\
DES & Smoothing & Both & Trend specification (records overstocking) \\
TES & Smoothing & Both & Full classical specification (weakest results) \\
ARIMA(1,1,0) & Statistical & Both & Linear memory \\
SARIMA(1,1,0)(0,1,1)$_7$ & Statistical & Store only & Seasonal challenger \\
LSTM (global) & Neural & Both & Pooled nonlinear learning \\
Croston & Intermittent & M5 only & Sparse-data specialist \\
SBA & Intermittent & M5 only & Debiased specialist \\
TSB & Intermittent & M5 only & Adaptive specialist \\
\bottomrule
\end{tabular}
\end{table}
